\documentclass[tikz,border=8pt]{standalone}
\usepackage[T1]{fontenc}
\usepackage{lmodern}
% mqttsystem-figure-system.tex
% Canonical TikZ visual system for MQTTSuite technical figures.
%
% Usage:
%   \documentclass[tikz,border=10pt]{standalone}
%   % mqttsystem-figure-system.tex
% Canonical TikZ visual system for MQTTSuite technical figures.
%
% Usage:
%   \documentclass[tikz,border=10pt]{standalone}
%   % mqttsystem-figure-system.tex
% Canonical TikZ visual system for MQTTSuite technical figures.
%
% Usage:
%   \documentclass[tikz,border=10pt]{standalone}
%   \input{mqttsystem-figure-system.tex}
%   \begin{document}
%   \begin{tikzpicture}[mqtt desktop]
%     ...
%   \end{tikzpicture}
%   \end{document}
%
% Canonical source rule:
% - figure .tex sources + this file are authoritative;
% - generated SVGs are publication output and are never hand-edited;
% - real runtime/UI evidence remains real raster capture.

\usetikzlibrary{
  arrows.meta,
  backgrounds,
  calc,
  fit,
  matrix,
  positioning,
  shapes.geometric,
  shapes.misc,
  shadows.blur
}

% ---------------------------------------------------------------------------
% Canonical palette
% ---------------------------------------------------------------------------

\definecolor{MQText}{HTML}{0F172A}
\definecolor{MQTextMuted}{HTML}{475569}
\definecolor{MQLine}{HTML}{475569}
\definecolor{MQBorder}{HTML}{CBD5E1}

\definecolor{MQSurface}{HTML}{F8FAFC}
\definecolor{MQSurfaceAlt}{HTML}{F1F5F9}
\definecolor{MQWhite}{HTML}{FFFFFF}

\definecolor{MQDecisionFill}{HTML}{FFF7ED}
\definecolor{MQDecisionLine}{HTML}{D97706}

\definecolor{MQSuccessFill}{HTML}{ECFDF5}
\definecolor{MQSuccessLine}{HTML}{10B981}

\definecolor{MQErrorFill}{HTML}{FEF2F2}
\definecolor{MQErrorLine}{HTML}{EF4444}

\definecolor{MQControlFill}{HTML}{EFF6FF}
\definecolor{MQControlLine}{HTML}{3B82F6}

\definecolor{MQWarningFill}{HTML}{FFF7ED}
\definecolor{MQWarningLine}{HTML}{F59E0B}

\definecolor{MQCrossLane}{HTML}{2563EB}

\definecolor{MQDataPlaneFill}{HTML}{EFF6FF}
\definecolor{MQDataPlaneLine}{HTML}{2563EB}

\definecolor{MQAdminPlaneFill}{HTML}{F8FAFC}
\definecolor{MQAdminPlaneLine}{HTML}{64748B}

\definecolor{MQTrustFill}{HTML}{FFF7ED}
\definecolor{MQTrustLine}{HTML}{D97706}

\definecolor{MQEvidenceRuntimeFill}{HTML}{ECFDF5}
\definecolor{MQEvidenceRuntimeLine}{HTML}{10B981}

\definecolor{MQEvidenceSourceFill}{HTML}{EFF6FF}
\definecolor{MQEvidenceSourceLine}{HTML}{3B82F6}

% ---------------------------------------------------------------------------
% Canonical radius system
% ---------------------------------------------------------------------------
%
% Small informational elements stay subtle; ordinary nodes are moderately
% rounded; structural containers are more rounded. Decision diamonds and
% database cylinders do not use these radii.

\newcommand{\MQRadiusXS}{1.2pt}
\newcommand{\MQRadiusS}{1.8pt}
\newcommand{\MQRadiusM}{2.8pt}
\newcommand{\MQRadiusL}{5.0pt}
\newcommand{\MQRadiusXL}{7.0pt}

% ---------------------------------------------------------------------------
% Canonical shading / elevation system
% ---------------------------------------------------------------------------
%
% Flat rendering is the default. Shadows are opt-in and reserved for semantic
% emphasis, never general decoration.

\newcommand{\MQShadowXSubtle}{0.7pt}
\newcommand{\MQShadowYSubtle}{-0.7pt}
\newcommand{\MQShadowBlurSubtle}{1.6pt}
\newcommand{\MQShadowOpacitySubtle}{0.12}

\newcommand{\MQShadowXRaised}{1.0pt}
\newcommand{\MQShadowYRaised}{-1.0pt}
\newcommand{\MQShadowBlurRaised}{2.4pt}
\newcommand{\MQShadowOpacityRaised}{0.16}

\tikzset{
  mqtt elevated none/.style={},
  mqtt elevated subtle/.style={
    blur shadow={
      shadow xshift=\MQShadowXSubtle,
      shadow yshift=\MQShadowYSubtle,
      shadow blur radius=\MQShadowBlurSubtle,
      opacity=\MQShadowOpacitySubtle
    }
  },
  mqtt elevated raised/.style={
    blur shadow={
      shadow xshift=\MQShadowXRaised,
      shadow yshift=\MQShadowYRaised,
      shadow blur radius=\MQShadowBlurRaised,
      opacity=\MQShadowOpacityRaised
    }
  }
}

% ---------------------------------------------------------------------------
% Canonical spacing system
% ---------------------------------------------------------------------------
%
% Base grid: 1 spacing unit = 2 mm. Ordinary figure spacing should use these
% semantic constants rather than ad-hoc gaps wherever practical. Exact geometry
% remains valid when a figure needs deterministic alignment, but should still
% follow this grid.

\newcommand{\MQU}{2mm}

% Generic scale
\newcommand{\MQGapXS}{2\MQU}      % 4 mm
\newcommand{\MQGapS}{3\MQU}       % 6 mm
\newcommand{\MQGapM}{5\MQU}       % 10 mm
\newcommand{\MQGapL}{7\MQU}       % 14 mm
\newcommand{\MQGapXL}{10\MQU}     % 20 mm

% Node interior
\newcommand{\MQNodePadX}{2.5\MQU} % 5 mm
\newcommand{\MQNodePadY}{2\MQU}   % 4 mm

% Sequential node spacing
\newcommand{\MQNodeGapTight}{3\MQU} % 6 mm
\newcommand{\MQNodeGap}{5\MQU}      % 10 mm
\newcommand{\MQNodeGapLoose}{7\MQU} % 14 mm

% Decision / branch geometry
\newcommand{\MQBranchGap}{7\MQU}    % 14 mm
\newcommand{\MQBranchRun}{5\MQU}    % 10 mm orthogonal run before turn

% Lanes and semantic containers
\newcommand{\MQLaneGap}{12\MQU}     % 24 mm between major lanes
\newcommand{\MQLanePadX}{4\MQU}     % 8 mm
\newcommand{\MQLanePadY}{4\MQU}     % 8 mm
\newcommand{\MQGroupPad}{3\MQU}     % 6 mm
\newcommand{\MQTrustPad}{4\MQU}     % 8 mm

% Titles and labels
\newcommand{\MQTitleGap}{2\MQU}     % 4 mm
\newcommand{\MQLabelGap}{1\MQU}     % 2 mm
\newcommand{\MQCrossLaneGap}{4\MQU} % 8 mm reserved clearance

% Figure safe area
\newcommand{\MQFigureSafeX}{5\MQU}  % 10 mm
\newcommand{\MQFigureSafeY}{5\MQU}  % 10 mm

% Compatibility aliases
\newcommand{\MQDesktopSafe}{\MQFigureSafeX}
\newcommand{\MQMobileSafe}{4\MQU}   % 8 mm

% Mobile art direction: mobile is not a scaled desktop composition.
\newcommand{\MQMobileNodeGap}{6\MQU}   % 12 mm
\newcommand{\MQMobileBranchGap}{8\MQU} % 16 mm
\newcommand{\MQMobileLaneGap}{10\MQU}  % 20 mm
\newcommand{\MQMobileGroupPad}{4\MQU}  % 8 mm

% ---------------------------------------------------------------------------
% Typography
% ---------------------------------------------------------------------------

\tikzset{
  mqtt text/.style={
    font=\sffamily\small,
    text=MQText,
    align=center
  },
  mqtt text left/.style={
    mqtt text,
    align=left
  },
  mqtt title/.style={
    font=\sffamily\bfseries\large,
    text=MQText,
    align=left
  },
  mqtt section title/.style={
    font=\sffamily\bfseries\small,
    text=MQText,
    align=left,
    anchor=west
  },
  mqtt annotation/.style={
    font=\sffamily\scriptsize,
    text=MQTextMuted,
    align=center
  },
  mqtt code/.style={
    font=\ttfamily\small,
    text=MQText
  },
  mqtt code small/.style={
    font=\ttfamily\scriptsize,
    text=MQText
  }
}

% ---------------------------------------------------------------------------
% Figure presets
% ---------------------------------------------------------------------------
%
% These establish common defaults only. Desktop/mobile variants may and often
% should use different geometry.

\tikzset{
  mqtt desktop/.style={
    x=1cm,
    y=1cm,
    every node/.append style={transform shape=false}
  },
  mqtt mobile/.style={
    x=1cm,
    y=1cm,
    every node/.append style={
      transform shape=false,
      font=\sffamily\normalsize
    }
  }
}

% ---------------------------------------------------------------------------
% Core node families
% ---------------------------------------------------------------------------

\tikzset{
  mqtt process/.style={
    mqtt text,
    draw=MQLine,
    fill=MQWhite,
    rounded corners=\MQRadiusM,
    line width=.75pt,
    minimum height=10mm,
    minimum width=25mm,
    inner xsep=\MQNodePadX,
    inner ysep=\MQNodePadY
  },
  mqtt application/.style={
    mqtt process,
    minimum width=31mm,
    minimum height=13mm,
    line width=.9pt
  },
  mqtt external/.style={
    mqtt process,
    fill=MQSurfaceAlt,
    draw=MQBorder
  },
  mqtt state/.style={
    mqtt process,
    fill=MQSurface
  },
  mqtt decision/.style={
    mqtt text,
    diamond,
    aspect=1.75,
    draw=MQDecisionLine,
    fill=MQDecisionFill,
    line width=.8pt,
    inner sep=1.2pt,
    minimum width=27mm,
    minimum height=14mm
  },
  mqtt success/.style={
    mqtt process,
    fill=MQSuccessFill,
    draw=MQSuccessLine
  },
  mqtt error/.style={
    mqtt process,
    fill=MQErrorFill,
    draw=MQErrorLine
  },
  mqtt control/.style={
    mqtt process,
    fill=MQControlFill,
    draw=MQControlLine
  },
  mqtt warning/.style={
    mqtt process,
    fill=MQWarningFill,
    draw=MQWarningLine
  },
  mqtt note/.style={
    mqtt text,
    draw=MQBorder,
    fill=MQSurfaceAlt,
    rounded corners=\MQRadiusS,
    line width=.65pt,
    font=\sffamily\scriptsize,
    minimum height=8mm,
    inner xsep=\MQNodePadX,
    inner ysep=\MQNodePadY
  },
  mqtt config/.style={
    mqtt process,
    fill=MQSurfaceAlt,
    draw=MQBorder,
    rounded corners=\MQRadiusS
  },
  mqtt data/.style={
    mqtt process,
    fill=MQWhite,
    draw=MQBorder
  },
  mqtt database/.style={
    mqtt text,
    cylinder,
    shape border rotate=90,
    aspect=.25,
    draw=MQLine,
    fill=MQWhite,
    line width=.75pt,
    minimum width=24mm,
    minimum height=13mm,
    inner xsep=\MQNodePadX,
    inner ysep=\MQNodePadY
  }
}

% ---------------------------------------------------------------------------
% Semantic containers
% ---------------------------------------------------------------------------

\tikzset{
  mqtt lane/.style={
    draw=MQBorder,
    fill=MQSurface,
    rounded corners=\MQRadiusL,
    line width=.8pt,
    inner xsep=\MQLanePadX,
    inner ysep=\MQLanePadY
  },
  mqtt group/.style={
    draw=MQBorder,
    fill=MQWhite,
    rounded corners=\MQRadiusL,
    line width=.75pt,
    inner sep=\MQGroupPad
  },
  mqtt data plane/.style={
    draw=MQDataPlaneLine,
    fill=MQDataPlaneFill,
    rounded corners=\MQRadiusL,
    line width=.9pt,
    inner sep=\MQGroupPad
  },
  mqtt admin plane/.style={
    draw=MQAdminPlaneLine,
    fill=MQAdminPlaneFill,
    rounded corners=\MQRadiusL,
    line width=.9pt,
    inner sep=\MQGroupPad
  },
  mqtt trust boundary/.style={
    draw=MQTrustLine,
    fill=none,
    rounded corners=\MQRadiusXL,
    line width=1.0pt,
    dashed,
    inner sep=\MQTrustPad
  }
}

% ---------------------------------------------------------------------------
% Arrow semantics
% ---------------------------------------------------------------------------

\tikzset{
  mqtt flow/.style={
    -{Stealth[length=2.2mm,width=1.5mm]},
    draw=MQLine,
    line width=.8pt
  },
  mqtt data flow/.style={
    -{Stealth[length=2.2mm,width=1.5mm]},
    draw=MQDataPlaneLine,
    line width=.9pt
  },
  mqtt control flow/.style={
    -{Stealth[length=2.2mm,width=1.5mm]},
    draw=MQControlLine,
    line width=.9pt
  },
  mqtt observation/.style={
    -{Stealth[length=2.0mm,width=1.4mm]},
    draw=MQLine,
    line width=.75pt,
    dashed
  },
  mqtt handoff/.style={
    -{Stealth[length=2.3mm,width=1.6mm]},
    draw=MQCrossLane,
    line width=1.0pt
  },
  mqtt association/.style={
    draw=MQBorder,
    line width=.7pt
  },
  mqtt dependency/.style={
    -{Stealth[length=2.0mm,width=1.4mm]},
    draw=MQLine,
    line width=.7pt,
    dashed
  }
}

% ---------------------------------------------------------------------------
% Labels and evidence
% ---------------------------------------------------------------------------

\tikzset{
  mqtt branch label/.style={
    font=\sffamily\scriptsize\bfseries,
    text=MQTextMuted,
    fill=MQWhite,
    inner sep=1.4pt
  },
  mqtt edge label/.style={
    font=\sffamily\scriptsize,
    text=MQTextMuted,
    fill=MQWhite,
    inner sep=1.4pt
  },
  mqtt topic label/.style={
    font=\ttfamily\scriptsize,
    text=MQText,
    fill=MQWhite,
    inner sep=1.4pt
  },
  mqtt route label/.style={
    font=\ttfamily\scriptsize,
    text=MQText,
    fill=MQWhite,
    inner sep=1.4pt
  },
  mqtt evidence runtime/.style={
    font=\sffamily\scriptsize\bfseries,
    text=MQSuccessLine,
    draw=MQEvidenceRuntimeLine,
    fill=MQEvidenceRuntimeFill,
    rounded corners=\MQRadiusXS,
    line width=.6pt,
    inner xsep=4pt,
    inner ysep=2pt
  },
  mqtt evidence source/.style={
    font=\sffamily\scriptsize\bfseries,
    text=MQControlLine,
    draw=MQEvidenceSourceLine,
    fill=MQEvidenceSourceFill,
    rounded corners=\MQRadiusXS,
    line width=.6pt,
    inner xsep=4pt,
    inner ysep=2pt
  }
}

% ---------------------------------------------------------------------------
% Canonical publication rules
% ---------------------------------------------------------------------------
%
% 1. Figure .tex + this style file are canonical technical sources.
% 2. Generated SVGs are never hand-edited.
% 3. SVG is the publication format for technical diagrams.
% 4. PNG is reserved for genuine runtime/UI evidence or raster-only input.
% 5. Desktop/mobile may share semantics but use independently art-directed geometry.
% 6. No arrow exists without real directional semantics.
% 7. Dashed arrows mean observation/secondary relation, not primary flow.
% 8. Containers denote real process/runtime/trust/configuration/lifecycle boundaries.
% 9. Color never carries the only semantic distinction.
% 10. Data plane, admin/control plane, trust boundaries and evidence classes use
%     the canonical semantic styles above.
% 11. Prefer explicit anchors/ports and orthogonal routing where clarity improves.
% 12. Preserve readable text size at actual GitHub desktop/mobile rendering widths.
% 13. Ordinary spacing follows the semantic spacing system above; avoid ad-hoc gaps.
% 14. Rounded corners follow the semantic radius scale above; avoid one-off radii.
% 15. Flat rendering is the default; elevation/shadows are explicit and semantically
%     justified.

%   \begin{document}
%   \begin{tikzpicture}[mqtt desktop]
%     ...
%   \end{tikzpicture}
%   \end{document}
%
% Canonical source rule:
% - figure .tex sources + this file are authoritative;
% - generated SVGs are publication output and are never hand-edited;
% - real runtime/UI evidence remains real raster capture.

\usetikzlibrary{
  arrows.meta,
  backgrounds,
  calc,
  fit,
  matrix,
  positioning,
  shapes.geometric,
  shapes.misc,
  shadows.blur
}

% ---------------------------------------------------------------------------
% Canonical palette
% ---------------------------------------------------------------------------

\definecolor{MQText}{HTML}{0F172A}
\definecolor{MQTextMuted}{HTML}{475569}
\definecolor{MQLine}{HTML}{475569}
\definecolor{MQBorder}{HTML}{CBD5E1}

\definecolor{MQSurface}{HTML}{F8FAFC}
\definecolor{MQSurfaceAlt}{HTML}{F1F5F9}
\definecolor{MQWhite}{HTML}{FFFFFF}

\definecolor{MQDecisionFill}{HTML}{FFF7ED}
\definecolor{MQDecisionLine}{HTML}{D97706}

\definecolor{MQSuccessFill}{HTML}{ECFDF5}
\definecolor{MQSuccessLine}{HTML}{10B981}

\definecolor{MQErrorFill}{HTML}{FEF2F2}
\definecolor{MQErrorLine}{HTML}{EF4444}

\definecolor{MQControlFill}{HTML}{EFF6FF}
\definecolor{MQControlLine}{HTML}{3B82F6}

\definecolor{MQWarningFill}{HTML}{FFF7ED}
\definecolor{MQWarningLine}{HTML}{F59E0B}

\definecolor{MQCrossLane}{HTML}{2563EB}

\definecolor{MQDataPlaneFill}{HTML}{EFF6FF}
\definecolor{MQDataPlaneLine}{HTML}{2563EB}

\definecolor{MQAdminPlaneFill}{HTML}{F8FAFC}
\definecolor{MQAdminPlaneLine}{HTML}{64748B}

\definecolor{MQTrustFill}{HTML}{FFF7ED}
\definecolor{MQTrustLine}{HTML}{D97706}

\definecolor{MQEvidenceRuntimeFill}{HTML}{ECFDF5}
\definecolor{MQEvidenceRuntimeLine}{HTML}{10B981}

\definecolor{MQEvidenceSourceFill}{HTML}{EFF6FF}
\definecolor{MQEvidenceSourceLine}{HTML}{3B82F6}

% ---------------------------------------------------------------------------
% Canonical radius system
% ---------------------------------------------------------------------------
%
% Small informational elements stay subtle; ordinary nodes are moderately
% rounded; structural containers are more rounded. Decision diamonds and
% database cylinders do not use these radii.

\newcommand{\MQRadiusXS}{1.2pt}
\newcommand{\MQRadiusS}{1.8pt}
\newcommand{\MQRadiusM}{2.8pt}
\newcommand{\MQRadiusL}{5.0pt}
\newcommand{\MQRadiusXL}{7.0pt}

% ---------------------------------------------------------------------------
% Canonical shading / elevation system
% ---------------------------------------------------------------------------
%
% Flat rendering is the default. Shadows are opt-in and reserved for semantic
% emphasis, never general decoration.

\newcommand{\MQShadowXSubtle}{0.7pt}
\newcommand{\MQShadowYSubtle}{-0.7pt}
\newcommand{\MQShadowBlurSubtle}{1.6pt}
\newcommand{\MQShadowOpacitySubtle}{0.12}

\newcommand{\MQShadowXRaised}{1.0pt}
\newcommand{\MQShadowYRaised}{-1.0pt}
\newcommand{\MQShadowBlurRaised}{2.4pt}
\newcommand{\MQShadowOpacityRaised}{0.16}

\tikzset{
  mqtt elevated none/.style={},
  mqtt elevated subtle/.style={
    blur shadow={
      shadow xshift=\MQShadowXSubtle,
      shadow yshift=\MQShadowYSubtle,
      shadow blur radius=\MQShadowBlurSubtle,
      opacity=\MQShadowOpacitySubtle
    }
  },
  mqtt elevated raised/.style={
    blur shadow={
      shadow xshift=\MQShadowXRaised,
      shadow yshift=\MQShadowYRaised,
      shadow blur radius=\MQShadowBlurRaised,
      opacity=\MQShadowOpacityRaised
    }
  }
}

% ---------------------------------------------------------------------------
% Canonical spacing system
% ---------------------------------------------------------------------------
%
% Base grid: 1 spacing unit = 2 mm. Ordinary figure spacing should use these
% semantic constants rather than ad-hoc gaps wherever practical. Exact geometry
% remains valid when a figure needs deterministic alignment, but should still
% follow this grid.

\newcommand{\MQU}{2mm}

% Generic scale
\newcommand{\MQGapXS}{2\MQU}      % 4 mm
\newcommand{\MQGapS}{3\MQU}       % 6 mm
\newcommand{\MQGapM}{5\MQU}       % 10 mm
\newcommand{\MQGapL}{7\MQU}       % 14 mm
\newcommand{\MQGapXL}{10\MQU}     % 20 mm

% Node interior
\newcommand{\MQNodePadX}{2.5\MQU} % 5 mm
\newcommand{\MQNodePadY}{2\MQU}   % 4 mm

% Sequential node spacing
\newcommand{\MQNodeGapTight}{3\MQU} % 6 mm
\newcommand{\MQNodeGap}{5\MQU}      % 10 mm
\newcommand{\MQNodeGapLoose}{7\MQU} % 14 mm

% Decision / branch geometry
\newcommand{\MQBranchGap}{7\MQU}    % 14 mm
\newcommand{\MQBranchRun}{5\MQU}    % 10 mm orthogonal run before turn

% Lanes and semantic containers
\newcommand{\MQLaneGap}{12\MQU}     % 24 mm between major lanes
\newcommand{\MQLanePadX}{4\MQU}     % 8 mm
\newcommand{\MQLanePadY}{4\MQU}     % 8 mm
\newcommand{\MQGroupPad}{3\MQU}     % 6 mm
\newcommand{\MQTrustPad}{4\MQU}     % 8 mm

% Titles and labels
\newcommand{\MQTitleGap}{2\MQU}     % 4 mm
\newcommand{\MQLabelGap}{1\MQU}     % 2 mm
\newcommand{\MQCrossLaneGap}{4\MQU} % 8 mm reserved clearance

% Figure safe area
\newcommand{\MQFigureSafeX}{5\MQU}  % 10 mm
\newcommand{\MQFigureSafeY}{5\MQU}  % 10 mm

% Compatibility aliases
\newcommand{\MQDesktopSafe}{\MQFigureSafeX}
\newcommand{\MQMobileSafe}{4\MQU}   % 8 mm

% Mobile art direction: mobile is not a scaled desktop composition.
\newcommand{\MQMobileNodeGap}{6\MQU}   % 12 mm
\newcommand{\MQMobileBranchGap}{8\MQU} % 16 mm
\newcommand{\MQMobileLaneGap}{10\MQU}  % 20 mm
\newcommand{\MQMobileGroupPad}{4\MQU}  % 8 mm

% ---------------------------------------------------------------------------
% Typography
% ---------------------------------------------------------------------------

\tikzset{
  mqtt text/.style={
    font=\sffamily\small,
    text=MQText,
    align=center
  },
  mqtt text left/.style={
    mqtt text,
    align=left
  },
  mqtt title/.style={
    font=\sffamily\bfseries\large,
    text=MQText,
    align=left
  },
  mqtt section title/.style={
    font=\sffamily\bfseries\small,
    text=MQText,
    align=left,
    anchor=west
  },
  mqtt annotation/.style={
    font=\sffamily\scriptsize,
    text=MQTextMuted,
    align=center
  },
  mqtt code/.style={
    font=\ttfamily\small,
    text=MQText
  },
  mqtt code small/.style={
    font=\ttfamily\scriptsize,
    text=MQText
  }
}

% ---------------------------------------------------------------------------
% Figure presets
% ---------------------------------------------------------------------------
%
% These establish common defaults only. Desktop/mobile variants may and often
% should use different geometry.

\tikzset{
  mqtt desktop/.style={
    x=1cm,
    y=1cm,
    every node/.append style={transform shape=false}
  },
  mqtt mobile/.style={
    x=1cm,
    y=1cm,
    every node/.append style={
      transform shape=false,
      font=\sffamily\normalsize
    }
  }
}

% ---------------------------------------------------------------------------
% Core node families
% ---------------------------------------------------------------------------

\tikzset{
  mqtt process/.style={
    mqtt text,
    draw=MQLine,
    fill=MQWhite,
    rounded corners=\MQRadiusM,
    line width=.75pt,
    minimum height=10mm,
    minimum width=25mm,
    inner xsep=\MQNodePadX,
    inner ysep=\MQNodePadY
  },
  mqtt application/.style={
    mqtt process,
    minimum width=31mm,
    minimum height=13mm,
    line width=.9pt
  },
  mqtt external/.style={
    mqtt process,
    fill=MQSurfaceAlt,
    draw=MQBorder
  },
  mqtt state/.style={
    mqtt process,
    fill=MQSurface
  },
  mqtt decision/.style={
    mqtt text,
    diamond,
    aspect=1.75,
    draw=MQDecisionLine,
    fill=MQDecisionFill,
    line width=.8pt,
    inner sep=1.2pt,
    minimum width=27mm,
    minimum height=14mm
  },
  mqtt success/.style={
    mqtt process,
    fill=MQSuccessFill,
    draw=MQSuccessLine
  },
  mqtt error/.style={
    mqtt process,
    fill=MQErrorFill,
    draw=MQErrorLine
  },
  mqtt control/.style={
    mqtt process,
    fill=MQControlFill,
    draw=MQControlLine
  },
  mqtt warning/.style={
    mqtt process,
    fill=MQWarningFill,
    draw=MQWarningLine
  },
  mqtt note/.style={
    mqtt text,
    draw=MQBorder,
    fill=MQSurfaceAlt,
    rounded corners=\MQRadiusS,
    line width=.65pt,
    font=\sffamily\scriptsize,
    minimum height=8mm,
    inner xsep=\MQNodePadX,
    inner ysep=\MQNodePadY
  },
  mqtt config/.style={
    mqtt process,
    fill=MQSurfaceAlt,
    draw=MQBorder,
    rounded corners=\MQRadiusS
  },
  mqtt data/.style={
    mqtt process,
    fill=MQWhite,
    draw=MQBorder
  },
  mqtt database/.style={
    mqtt text,
    cylinder,
    shape border rotate=90,
    aspect=.25,
    draw=MQLine,
    fill=MQWhite,
    line width=.75pt,
    minimum width=24mm,
    minimum height=13mm,
    inner xsep=\MQNodePadX,
    inner ysep=\MQNodePadY
  }
}

% ---------------------------------------------------------------------------
% Semantic containers
% ---------------------------------------------------------------------------

\tikzset{
  mqtt lane/.style={
    draw=MQBorder,
    fill=MQSurface,
    rounded corners=\MQRadiusL,
    line width=.8pt,
    inner xsep=\MQLanePadX,
    inner ysep=\MQLanePadY
  },
  mqtt group/.style={
    draw=MQBorder,
    fill=MQWhite,
    rounded corners=\MQRadiusL,
    line width=.75pt,
    inner sep=\MQGroupPad
  },
  mqtt data plane/.style={
    draw=MQDataPlaneLine,
    fill=MQDataPlaneFill,
    rounded corners=\MQRadiusL,
    line width=.9pt,
    inner sep=\MQGroupPad
  },
  mqtt admin plane/.style={
    draw=MQAdminPlaneLine,
    fill=MQAdminPlaneFill,
    rounded corners=\MQRadiusL,
    line width=.9pt,
    inner sep=\MQGroupPad
  },
  mqtt trust boundary/.style={
    draw=MQTrustLine,
    fill=none,
    rounded corners=\MQRadiusXL,
    line width=1.0pt,
    dashed,
    inner sep=\MQTrustPad
  }
}

% ---------------------------------------------------------------------------
% Arrow semantics
% ---------------------------------------------------------------------------

\tikzset{
  mqtt flow/.style={
    -{Stealth[length=2.2mm,width=1.5mm]},
    draw=MQLine,
    line width=.8pt
  },
  mqtt data flow/.style={
    -{Stealth[length=2.2mm,width=1.5mm]},
    draw=MQDataPlaneLine,
    line width=.9pt
  },
  mqtt control flow/.style={
    -{Stealth[length=2.2mm,width=1.5mm]},
    draw=MQControlLine,
    line width=.9pt
  },
  mqtt observation/.style={
    -{Stealth[length=2.0mm,width=1.4mm]},
    draw=MQLine,
    line width=.75pt,
    dashed
  },
  mqtt handoff/.style={
    -{Stealth[length=2.3mm,width=1.6mm]},
    draw=MQCrossLane,
    line width=1.0pt
  },
  mqtt association/.style={
    draw=MQBorder,
    line width=.7pt
  },
  mqtt dependency/.style={
    -{Stealth[length=2.0mm,width=1.4mm]},
    draw=MQLine,
    line width=.7pt,
    dashed
  }
}

% ---------------------------------------------------------------------------
% Labels and evidence
% ---------------------------------------------------------------------------

\tikzset{
  mqtt branch label/.style={
    font=\sffamily\scriptsize\bfseries,
    text=MQTextMuted,
    fill=MQWhite,
    inner sep=1.4pt
  },
  mqtt edge label/.style={
    font=\sffamily\scriptsize,
    text=MQTextMuted,
    fill=MQWhite,
    inner sep=1.4pt
  },
  mqtt topic label/.style={
    font=\ttfamily\scriptsize,
    text=MQText,
    fill=MQWhite,
    inner sep=1.4pt
  },
  mqtt route label/.style={
    font=\ttfamily\scriptsize,
    text=MQText,
    fill=MQWhite,
    inner sep=1.4pt
  },
  mqtt evidence runtime/.style={
    font=\sffamily\scriptsize\bfseries,
    text=MQSuccessLine,
    draw=MQEvidenceRuntimeLine,
    fill=MQEvidenceRuntimeFill,
    rounded corners=\MQRadiusXS,
    line width=.6pt,
    inner xsep=4pt,
    inner ysep=2pt
  },
  mqtt evidence source/.style={
    font=\sffamily\scriptsize\bfseries,
    text=MQControlLine,
    draw=MQEvidenceSourceLine,
    fill=MQEvidenceSourceFill,
    rounded corners=\MQRadiusXS,
    line width=.6pt,
    inner xsep=4pt,
    inner ysep=2pt
  }
}

% ---------------------------------------------------------------------------
% Canonical publication rules
% ---------------------------------------------------------------------------
%
% 1. Figure .tex + this style file are canonical technical sources.
% 2. Generated SVGs are never hand-edited.
% 3. SVG is the publication format for technical diagrams.
% 4. PNG is reserved for genuine runtime/UI evidence or raster-only input.
% 5. Desktop/mobile may share semantics but use independently art-directed geometry.
% 6. No arrow exists without real directional semantics.
% 7. Dashed arrows mean observation/secondary relation, not primary flow.
% 8. Containers denote real process/runtime/trust/configuration/lifecycle boundaries.
% 9. Color never carries the only semantic distinction.
% 10. Data plane, admin/control plane, trust boundaries and evidence classes use
%     the canonical semantic styles above.
% 11. Prefer explicit anchors/ports and orthogonal routing where clarity improves.
% 12. Preserve readable text size at actual GitHub desktop/mobile rendering widths.
% 13. Ordinary spacing follows the semantic spacing system above; avoid ad-hoc gaps.
% 14. Rounded corners follow the semantic radius scale above; avoid one-off radii.
% 15. Flat rendering is the default; elevation/shadows are explicit and semantically
%     justified.

%   \begin{document}
%   \begin{tikzpicture}[mqtt desktop]
%     ...
%   \end{tikzpicture}
%   \end{document}
%
% Canonical source rule:
% - figure .tex sources + this file are authoritative;
% - generated SVGs are publication output and are never hand-edited;
% - real runtime/UI evidence remains real raster capture.

\usetikzlibrary{
  arrows.meta,
  backgrounds,
  calc,
  fit,
  matrix,
  positioning,
  shapes.geometric,
  shapes.misc,
  shadows.blur
}

% ---------------------------------------------------------------------------
% Canonical palette
% ---------------------------------------------------------------------------

\definecolor{MQText}{HTML}{0F172A}
\definecolor{MQTextMuted}{HTML}{475569}
\definecolor{MQLine}{HTML}{475569}
\definecolor{MQBorder}{HTML}{CBD5E1}

\definecolor{MQSurface}{HTML}{F8FAFC}
\definecolor{MQSurfaceAlt}{HTML}{F1F5F9}
\definecolor{MQWhite}{HTML}{FFFFFF}

\definecolor{MQDecisionFill}{HTML}{FFF7ED}
\definecolor{MQDecisionLine}{HTML}{D97706}

\definecolor{MQSuccessFill}{HTML}{ECFDF5}
\definecolor{MQSuccessLine}{HTML}{10B981}

\definecolor{MQErrorFill}{HTML}{FEF2F2}
\definecolor{MQErrorLine}{HTML}{EF4444}

\definecolor{MQControlFill}{HTML}{EFF6FF}
\definecolor{MQControlLine}{HTML}{3B82F6}

\definecolor{MQWarningFill}{HTML}{FFF7ED}
\definecolor{MQWarningLine}{HTML}{F59E0B}

\definecolor{MQCrossLane}{HTML}{2563EB}

\definecolor{MQDataPlaneFill}{HTML}{EFF6FF}
\definecolor{MQDataPlaneLine}{HTML}{2563EB}

\definecolor{MQAdminPlaneFill}{HTML}{F8FAFC}
\definecolor{MQAdminPlaneLine}{HTML}{64748B}

\definecolor{MQTrustFill}{HTML}{FFF7ED}
\definecolor{MQTrustLine}{HTML}{D97706}

\definecolor{MQEvidenceRuntimeFill}{HTML}{ECFDF5}
\definecolor{MQEvidenceRuntimeLine}{HTML}{10B981}

\definecolor{MQEvidenceSourceFill}{HTML}{EFF6FF}
\definecolor{MQEvidenceSourceLine}{HTML}{3B82F6}

% ---------------------------------------------------------------------------
% Canonical radius system
% ---------------------------------------------------------------------------
%
% Small informational elements stay subtle; ordinary nodes are moderately
% rounded; structural containers are more rounded. Decision diamonds and
% database cylinders do not use these radii.

\newcommand{\MQRadiusXS}{1.2pt}
\newcommand{\MQRadiusS}{1.8pt}
\newcommand{\MQRadiusM}{2.8pt}
\newcommand{\MQRadiusL}{5.0pt}
\newcommand{\MQRadiusXL}{7.0pt}

% ---------------------------------------------------------------------------
% Canonical shading / elevation system
% ---------------------------------------------------------------------------
%
% Flat rendering is the default. Shadows are opt-in and reserved for semantic
% emphasis, never general decoration.

\newcommand{\MQShadowXSubtle}{0.7pt}
\newcommand{\MQShadowYSubtle}{-0.7pt}
\newcommand{\MQShadowBlurSubtle}{1.6pt}
\newcommand{\MQShadowOpacitySubtle}{0.12}

\newcommand{\MQShadowXRaised}{1.0pt}
\newcommand{\MQShadowYRaised}{-1.0pt}
\newcommand{\MQShadowBlurRaised}{2.4pt}
\newcommand{\MQShadowOpacityRaised}{0.16}

\tikzset{
  mqtt elevated none/.style={},
  mqtt elevated subtle/.style={
    blur shadow={
      shadow xshift=\MQShadowXSubtle,
      shadow yshift=\MQShadowYSubtle,
      shadow blur radius=\MQShadowBlurSubtle,
      opacity=\MQShadowOpacitySubtle
    }
  },
  mqtt elevated raised/.style={
    blur shadow={
      shadow xshift=\MQShadowXRaised,
      shadow yshift=\MQShadowYRaised,
      shadow blur radius=\MQShadowBlurRaised,
      opacity=\MQShadowOpacityRaised
    }
  }
}

% ---------------------------------------------------------------------------
% Canonical spacing system
% ---------------------------------------------------------------------------
%
% Base grid: 1 spacing unit = 2 mm. Ordinary figure spacing should use these
% semantic constants rather than ad-hoc gaps wherever practical. Exact geometry
% remains valid when a figure needs deterministic alignment, but should still
% follow this grid.

\newcommand{\MQU}{2mm}

% Generic scale
\newcommand{\MQGapXS}{2\MQU}      % 4 mm
\newcommand{\MQGapS}{3\MQU}       % 6 mm
\newcommand{\MQGapM}{5\MQU}       % 10 mm
\newcommand{\MQGapL}{7\MQU}       % 14 mm
\newcommand{\MQGapXL}{10\MQU}     % 20 mm

% Node interior
\newcommand{\MQNodePadX}{2.5\MQU} % 5 mm
\newcommand{\MQNodePadY}{2\MQU}   % 4 mm

% Sequential node spacing
\newcommand{\MQNodeGapTight}{3\MQU} % 6 mm
\newcommand{\MQNodeGap}{5\MQU}      % 10 mm
\newcommand{\MQNodeGapLoose}{7\MQU} % 14 mm

% Decision / branch geometry
\newcommand{\MQBranchGap}{7\MQU}    % 14 mm
\newcommand{\MQBranchRun}{5\MQU}    % 10 mm orthogonal run before turn

% Lanes and semantic containers
\newcommand{\MQLaneGap}{12\MQU}     % 24 mm between major lanes
\newcommand{\MQLanePadX}{4\MQU}     % 8 mm
\newcommand{\MQLanePadY}{4\MQU}     % 8 mm
\newcommand{\MQGroupPad}{3\MQU}     % 6 mm
\newcommand{\MQTrustPad}{4\MQU}     % 8 mm

% Titles and labels
\newcommand{\MQTitleGap}{2\MQU}     % 4 mm
\newcommand{\MQLabelGap}{1\MQU}     % 2 mm
\newcommand{\MQCrossLaneGap}{4\MQU} % 8 mm reserved clearance

% Figure safe area
\newcommand{\MQFigureSafeX}{5\MQU}  % 10 mm
\newcommand{\MQFigureSafeY}{5\MQU}  % 10 mm

% Compatibility aliases
\newcommand{\MQDesktopSafe}{\MQFigureSafeX}
\newcommand{\MQMobileSafe}{4\MQU}   % 8 mm

% Mobile art direction: mobile is not a scaled desktop composition.
\newcommand{\MQMobileNodeGap}{6\MQU}   % 12 mm
\newcommand{\MQMobileBranchGap}{8\MQU} % 16 mm
\newcommand{\MQMobileLaneGap}{10\MQU}  % 20 mm
\newcommand{\MQMobileGroupPad}{4\MQU}  % 8 mm

% ---------------------------------------------------------------------------
% Typography
% ---------------------------------------------------------------------------

\tikzset{
  mqtt text/.style={
    font=\sffamily\small,
    text=MQText,
    align=center
  },
  mqtt text left/.style={
    mqtt text,
    align=left
  },
  mqtt title/.style={
    font=\sffamily\bfseries\large,
    text=MQText,
    align=left
  },
  mqtt section title/.style={
    font=\sffamily\bfseries\small,
    text=MQText,
    align=left,
    anchor=west
  },
  mqtt annotation/.style={
    font=\sffamily\scriptsize,
    text=MQTextMuted,
    align=center
  },
  mqtt code/.style={
    font=\ttfamily\small,
    text=MQText
  },
  mqtt code small/.style={
    font=\ttfamily\scriptsize,
    text=MQText
  }
}

% ---------------------------------------------------------------------------
% Figure presets
% ---------------------------------------------------------------------------
%
% These establish common defaults only. Desktop/mobile variants may and often
% should use different geometry.

\tikzset{
  mqtt desktop/.style={
    x=1cm,
    y=1cm,
    every node/.append style={transform shape=false}
  },
  mqtt mobile/.style={
    x=1cm,
    y=1cm,
    every node/.append style={
      transform shape=false,
      font=\sffamily\normalsize
    }
  }
}

% ---------------------------------------------------------------------------
% Core node families
% ---------------------------------------------------------------------------

\tikzset{
  mqtt process/.style={
    mqtt text,
    draw=MQLine,
    fill=MQWhite,
    rounded corners=\MQRadiusM,
    line width=.75pt,
    minimum height=10mm,
    minimum width=25mm,
    inner xsep=\MQNodePadX,
    inner ysep=\MQNodePadY
  },
  mqtt application/.style={
    mqtt process,
    minimum width=31mm,
    minimum height=13mm,
    line width=.9pt
  },
  mqtt external/.style={
    mqtt process,
    fill=MQSurfaceAlt,
    draw=MQBorder
  },
  mqtt state/.style={
    mqtt process,
    fill=MQSurface
  },
  mqtt decision/.style={
    mqtt text,
    diamond,
    aspect=1.75,
    draw=MQDecisionLine,
    fill=MQDecisionFill,
    line width=.8pt,
    inner sep=1.2pt,
    minimum width=27mm,
    minimum height=14mm
  },
  mqtt success/.style={
    mqtt process,
    fill=MQSuccessFill,
    draw=MQSuccessLine
  },
  mqtt error/.style={
    mqtt process,
    fill=MQErrorFill,
    draw=MQErrorLine
  },
  mqtt control/.style={
    mqtt process,
    fill=MQControlFill,
    draw=MQControlLine
  },
  mqtt warning/.style={
    mqtt process,
    fill=MQWarningFill,
    draw=MQWarningLine
  },
  mqtt note/.style={
    mqtt text,
    draw=MQBorder,
    fill=MQSurfaceAlt,
    rounded corners=\MQRadiusS,
    line width=.65pt,
    font=\sffamily\scriptsize,
    minimum height=8mm,
    inner xsep=\MQNodePadX,
    inner ysep=\MQNodePadY
  },
  mqtt config/.style={
    mqtt process,
    fill=MQSurfaceAlt,
    draw=MQBorder,
    rounded corners=\MQRadiusS
  },
  mqtt data/.style={
    mqtt process,
    fill=MQWhite,
    draw=MQBorder
  },
  mqtt database/.style={
    mqtt text,
    cylinder,
    shape border rotate=90,
    aspect=.25,
    draw=MQLine,
    fill=MQWhite,
    line width=.75pt,
    minimum width=24mm,
    minimum height=13mm,
    inner xsep=\MQNodePadX,
    inner ysep=\MQNodePadY
  }
}

% ---------------------------------------------------------------------------
% Semantic containers
% ---------------------------------------------------------------------------

\tikzset{
  mqtt lane/.style={
    draw=MQBorder,
    fill=MQSurface,
    rounded corners=\MQRadiusL,
    line width=.8pt,
    inner xsep=\MQLanePadX,
    inner ysep=\MQLanePadY
  },
  mqtt group/.style={
    draw=MQBorder,
    fill=MQWhite,
    rounded corners=\MQRadiusL,
    line width=.75pt,
    inner sep=\MQGroupPad
  },
  mqtt data plane/.style={
    draw=MQDataPlaneLine,
    fill=MQDataPlaneFill,
    rounded corners=\MQRadiusL,
    line width=.9pt,
    inner sep=\MQGroupPad
  },
  mqtt admin plane/.style={
    draw=MQAdminPlaneLine,
    fill=MQAdminPlaneFill,
    rounded corners=\MQRadiusL,
    line width=.9pt,
    inner sep=\MQGroupPad
  },
  mqtt trust boundary/.style={
    draw=MQTrustLine,
    fill=none,
    rounded corners=\MQRadiusXL,
    line width=1.0pt,
    dashed,
    inner sep=\MQTrustPad
  }
}

% ---------------------------------------------------------------------------
% Arrow semantics
% ---------------------------------------------------------------------------

\tikzset{
  mqtt flow/.style={
    -{Stealth[length=2.2mm,width=1.5mm]},
    draw=MQLine,
    line width=.8pt
  },
  mqtt data flow/.style={
    -{Stealth[length=2.2mm,width=1.5mm]},
    draw=MQDataPlaneLine,
    line width=.9pt
  },
  mqtt control flow/.style={
    -{Stealth[length=2.2mm,width=1.5mm]},
    draw=MQControlLine,
    line width=.9pt
  },
  mqtt observation/.style={
    -{Stealth[length=2.0mm,width=1.4mm]},
    draw=MQLine,
    line width=.75pt,
    dashed
  },
  mqtt handoff/.style={
    -{Stealth[length=2.3mm,width=1.6mm]},
    draw=MQCrossLane,
    line width=1.0pt
  },
  mqtt association/.style={
    draw=MQBorder,
    line width=.7pt
  },
  mqtt dependency/.style={
    -{Stealth[length=2.0mm,width=1.4mm]},
    draw=MQLine,
    line width=.7pt,
    dashed
  }
}

% ---------------------------------------------------------------------------
% Labels and evidence
% ---------------------------------------------------------------------------

\tikzset{
  mqtt branch label/.style={
    font=\sffamily\scriptsize\bfseries,
    text=MQTextMuted,
    fill=MQWhite,
    inner sep=1.4pt
  },
  mqtt edge label/.style={
    font=\sffamily\scriptsize,
    text=MQTextMuted,
    fill=MQWhite,
    inner sep=1.4pt
  },
  mqtt topic label/.style={
    font=\ttfamily\scriptsize,
    text=MQText,
    fill=MQWhite,
    inner sep=1.4pt
  },
  mqtt route label/.style={
    font=\ttfamily\scriptsize,
    text=MQText,
    fill=MQWhite,
    inner sep=1.4pt
  },
  mqtt evidence runtime/.style={
    font=\sffamily\scriptsize\bfseries,
    text=MQSuccessLine,
    draw=MQEvidenceRuntimeLine,
    fill=MQEvidenceRuntimeFill,
    rounded corners=\MQRadiusXS,
    line width=.6pt,
    inner xsep=4pt,
    inner ysep=2pt
  },
  mqtt evidence source/.style={
    font=\sffamily\scriptsize\bfseries,
    text=MQControlLine,
    draw=MQEvidenceSourceLine,
    fill=MQEvidenceSourceFill,
    rounded corners=\MQRadiusXS,
    line width=.6pt,
    inner xsep=4pt,
    inner ysep=2pt
  }
}

% ---------------------------------------------------------------------------
% Canonical publication rules
% ---------------------------------------------------------------------------
%
% 1. Figure .tex + this style file are canonical technical sources.
% 2. Generated SVGs are never hand-edited.
% 3. SVG is the publication format for technical diagrams.
% 4. PNG is reserved for genuine runtime/UI evidence or raster-only input.
% 5. Desktop/mobile may share semantics but use independently art-directed geometry.
% 6. No arrow exists without real directional semantics.
% 7. Dashed arrows mean observation/secondary relation, not primary flow.
% 8. Containers denote real process/runtime/trust/configuration/lifecycle boundaries.
% 9. Color never carries the only semantic distinction.
% 10. Data plane, admin/control plane, trust boundaries and evidence classes use
%     the canonical semantic styles above.
% 11. Prefer explicit anchors/ports and orthogonal routing where clarity improves.
% 12. Preserve readable text size at actual GitHub desktop/mobile rendering widths.
% 13. Ordinary spacing follows the semantic spacing system above; avoid ad-hoc gaps.
% 14. Rounded corners follow the semantic radius scale above; avoid one-off radii.
% 15. Flat rendering is the default; elevation/shadows are explicit and semantically
%     justified.


\begin{document}
\begin{tikzpicture}[mqtt desktop]

\tikzset{
  lifecycle node/.style={minimum width=30mm,text width=27mm},
  lifecycle decision/.style={minimum width=30mm,text width=25mm}
}

% ===========================================================================
% HTTP administration
% ===========================================================================

\node[mqtt process,lifecycle node] (h0) {Operator};

\node[mqtt process,lifecycle node,
      right=\MQHorizontalGap of h0] (h1)
  {PATCH\\\texttt{/api/bridge/config}};

\node[mqtt decision,lifecycle decision,
      right=\MQHorizontalGap of h1] (h2)
  {Restart\\in progress?};

\node[mqtt control,lifecycle node,
      right=\MQHorizontalGap of h2] (h3)
  {Apply JSON Patch\\validate + defaults\\stage config};

\node[mqtt success,lifecycle node,
      right=\MQHorizontalGap of h3] (h4)
  {HTTP 200\\patch accepted};

\node[mqtt error,lifecycle node,
      below=\MQBranchGapTight of h2] (h2e)
  {HTTP 409\\patch not applied};

\node[mqtt error,lifecycle node,
      below=\MQBranchGapTight of h3] (h3e)
  {HTTP 404\\patch / validation failed};

\draw[mqtt flow] (h0.east) -- (h1.west);
\draw[mqtt flow] (h1.east) -- (h2.west);
\draw[mqtt flow] (h2.east) -- node[mqtt branch label,above]{NO} (h3.west);
\draw[mqtt flow] (h3.east) -- (h4.west);
\draw[mqtt flow] (h2.south) -- node[mqtt branch label,right]{YES} (h2e.north);
\draw[mqtt flow] (h3.south) -- (h3e.north);

% ===========================================================================
% Runtime restart lifecycle
% Anchor rightmost node below HTTP result, then build leftward.
% ===========================================================================

\node[mqtt process,lifecycle node,
      below=44mm of h4] (r0)
  {\texttt{closeBridges()}};

\node[mqtt process,lifecycle node,
      left=\MQHorizontalGap of r0] (r1)
  {Complete teardown\\disconnect / terminate\\active flows\\
   {\scriptsize immediate when none remain}};

\node[mqtt control,lifecycle node,
      left=\MQHorizontalGap of r1] (r2)
  {\texttt{activateStaged()}\\+\ \texttt{startBridges()}\\persist + rebuild};

\node[mqtt process,lifecycle node,
      left=\MQHorizontalGap of r2] (r3)
  {Broker clients\\connecting};

\node[mqtt success,lifecycle node,
      left=\MQHorizontalGap of r3] (r4)
  {Configured bridges\\connected / disabled};

\draw[mqtt flow] (r0.west) -- (r1.east);
\draw[mqtt flow] (r1.west) -- (r2.east);
\draw[mqtt flow] (r2.west) -- (r3.east);
\draw[mqtt flow] (r3.west) -- (r4.east);

% ===========================================================================
% SSE observation
% Anchor rightmost node below closeBridges(), then use identical widths/gaps.
% ===========================================================================

\node[mqtt external,lifecycle node,
      below=24mm of r0] (s0)
  {SSE observer};

\node[mqtt process,lifecycle node,
      left=\MQHorizontalGap of s0] (s1)
  {Stopping / disconnect events};

\node[mqtt process,lifecycle node,
      left=\MQHorizontalGap of s1] (s2)
  {\texttt{bridges\_starting}\\replay reset};

\node[mqtt process,lifecycle node,
      left=\MQHorizontalGap of s2] (s3)
  {Startup / connect events};

\node[mqtt success,lifecycle node,
      left=\MQHorizontalGap of s3] (s4)
  {\texttt{bridges\_started}};

\draw[mqtt flow] (s0.west) -- (s1.east);
\draw[mqtt flow] (s1.west) -- (s2.east);
\draw[mqtt flow] (s2.west) -- (s3.east);
\draw[mqtt flow] (s3.west) -- (s4.east);

\node[mqtt annotation,anchor=west,text width=75mm,
      below=5mm of s4.south west] (s-note)
  {Replay covers the current restart cycle only; \texttt{Last-Event-ID} is ignored; keep-alive every 39 s.};

% ===========================================================================
% Cross-lane semantics — exact vertical alignment
% ===========================================================================

\draw[mqtt handoff]
  (h4.south)
  -- node[mqtt edge label,right,text=MQCrossLane]
     {restart continues asynchronously}
  (r0.north);

\draw[mqtt observation] (r1.south) -- (s1.north);
\draw[mqtt observation] (r2.south) -- (s2.north);
\draw[mqtt observation] (r3.south) -- (s3.north);
\draw[mqtt observation] (r4.south) -- (s4.north);

% ===========================================================================
% Semantic lane containers + self-contained light background
% ===========================================================================

\begin{pgfonlayer}{background}
  \node[mqtt lane,fit=(h0)(h4)(h2e)(h3e)] (lane-http) {};
  \node[mqtt lane,fit=(r4)(r0)] (lane-runtime) {};
  \node[mqtt lane,fit=(s4)(s0)(s-note)] (lane-sse) {};

  \node[
    fill=MQWhite,
    draw=none,
    rounded corners=\MQRadiusL,
    fit=(lane-http)(lane-runtime)(lane-sse),
    inner sep=\MQFigureSafeX
  ] (figure-bg) {};
\end{pgfonlayer}

\node[mqtt section title]
  at ($(lane-http.north west)+(\MQTitleGap,-\MQTitleGap)$)
  {HTTP administration};

\node[mqtt section title]
  at ($(lane-runtime.north west)+(\MQTitleGap,-\MQTitleGap)$)
  {Runtime restart lifecycle};

\node[mqtt section title]
  at ($(lane-sse.north west)+(\MQTitleGap,-\MQTitleGap)$)
  {SSE observation};

\node[mqtt annotation,anchor=east]
  at ($(lane-sse.north east)+(-\MQTitleGap,-\MQTitleGap)$)
  {dashed vertical arrows = runtime lifecycle events observed via SSE};

\end{tikzpicture}
\end{document}
